\FloatBarrier

\section{Frozen-latent spectrum predictor}
\label{sec:spectral-predictor}

Finally, we test whether memorization onset can be estimated from the frozen VAE representation before training the diffusion model. For each trained VAE we encode the same diverse 1k subset, form the diffused latent covariance spectrum, and use its eigenvalues to define a mode-absorption clock $P_d(s)$ over training step $s$. We then pass this clock through one shared pressure-to-memorization response curve, giving a predicted memorization percentage over training for each $\dlat$. Sec.~\ref{sec:theory-predictor} derives what such a clock should contain: the covariance spectrum labels only the $\dlat$ data-side modes of the feature correlation matrix, whereas memorization is carried by the $\nsamp-\dlat$ sample-specific modes, whose noise-averaged eigenvalue $g(q_\mu)/\nsamp$ depends on the spectrum only through $q_\mu$. The same response is used across latent widths within a dataset, so the test asks whether frozen latent geometry contains the right temporal ordering; Appendix~\ref{app:spectral-predictor-clean} defines each term and gives an $n$-dependence diagnostic in Figure~\ref{fig:app-spectral-by-n}.

\begin{figure*}[!t]
  \centering
  \includegraphics[width=0.98\textwidth,height=0.58\textheight,keepaspectratio]{figures/real_spectral_predictor_compact_2x2.pdf}
  \caption{\textbf{Spectral memorization predictor on real datasets.} Representative 2$\times$2 blocks are shown for CelebA and CIFAR-10; the full sweeps are in Appendix Figures~\ref{fig:app-spectral-celeba-small} and~\ref{fig:app-spectral-cifar-small}. Blue curves are five-seed empirical KNN memorization fractions with 95\% confidence bands. The dashed red curve maps the frozen-latent absorption clock through one shared sigmoid response from spectral pressure to memorization percentage; black circles mark observed iso-memorization hitting times, and red crosses mark predicted hitting times. The left edge of each sweep is representation-limited rather than a clean excess-buffer regime: below the empirical VAE identity threshold $\widehat d_{\rm ID}^{\rm VAE}(0.95)$, reconstructions do not preserve enough sample identity for pixel-space memorization to be detected reliably. Once the latent code is expressive enough, the frozen-latent spectrum captures the main temporal structure: memorization rises near the approximate intrinsic-dimensionality transition, then slows and decreases as increasing $\dlat$ adds excess latent buffer.}
  \label{fig:real-spectral-predictor}
  \label{fig:real-spectral-predictor-celeba}
  \label{fig:real-spectral-predictor-cifar}
\end{figure*}

Concretely, for encoded training latents $Z_d\in\mathbb R^{n\times d}$ and diffusion time $t=0.1$, we compute \[ M_t(d)=e^{-2t}Z_d^\top Z_d/n+(1-e^{-2t})I_d \] and let $\{\lambda_i(d)\}$ be its eigenvalues. The absorption clock is \[ P_d(s)= \frac{\sum_i w_i(d)\left(1-\exp[-\kappa\lambda_i(d)s]\right)^2} {\sum_i w_i(d)}, \] with $w_i=1$ for the primary plots. The predicted memorization fraction is an anchored sigmoid of this clock, $\widehat m_d(s)=[\sigma(a(P_d(s)-b))-\sigma(-ab)]/[1-\sigma(-ab)]$, using one set of $(\kappa,a,b)$ across the post-identity-threshold dimensions of a dataset. Iso-memorization times are then read off as $\hat\tau_q(d)=\inf\{s:\widehat m_d(s)\ge q\}$.

The complete spectral-predictor diagnostics and term definitions are in Appendix~\ref{app:spectral-predictor-clean}: aggregate overlays in Figure~\ref{fig:app-spectral-overlays}, full small multiples in Figures~\ref{fig:app-spectral-celeba-small} and \ref{fig:app-spectral-cifar-small}, predicted-vs-observed hitting-time summaries in Figure~\ref{fig:app-spectral-tau}, spectral pressure curves in Figure~\ref{fig:app-spectral-pressure}, and frozen-latent spectral features in Figures~\ref{fig:app-spectral-features} and \ref{fig:app-spectral-features-cifar}. The post-threshold dimensions, where the VAE preserves pixel-space identity, are the relevant test of the excess-buffer prediction; deviations at smaller $\dlat$ are expected because the representation itself suppresses measurable memorization.
